\section{Desarrollo}
Para el desarrollo de este trabajo se utilizará principalmente herramientas de Python, como lo son las librerías \textit{librosa} y \textit{pytorch}, las cuales ofrecen una amplia gama de funcionalidades para el procesamiento de audio y la extracción de características musicales. Estas librerías permiten realizar tareas como la lectura de archivos de audio, la generación de espectrogramas, la aplicación de la Transformada de Fourier y la extracción de características cromáticas, entre otras.
Como avance entregable se presentará el entrenamiento y procesamiento de datos con la implementación de un modelo de aprendizaje CNN 2‑D, el cual se entrenará utilizando un dataset de audio de guitarra acústica de \textit{Guitarset} con etiquetas de acordes. Donde se evaluará el desempeño del modelo utilizando métricas como exact‑match accuracy, top‑3 accuracy y macro‑F1, con el objetivo de alcanzar una precisión exacta mínima del 85 $\%$ en la identificación de acordes. Además, se documentará el proceso de desarrollo y se prepararán los scripts necesarios para la reproducción de los resultados obtenidos, así como también una pequeña interfaz de usuario para facilitar su uso.